\documentclass[11pt]{article}
\usepackage[margin=1in]{geometry}
\usepackage[T1]{fontenc}
\usepackage[utf8]{inputenc}
\usepackage{lmodern}
\usepackage{microtype}
\usepackage[table]{xcolor}
\usepackage{amsmath,amssymb,amsfonts,mathtools,bm}
\IfFileExists{physics.sty}{%
  \usepackage{physics}%
}{%
  % Portable subset used by this project when the optional physics package
  % is not installed in the local TeX distribution.
  \NewDocumentCommand{\pdv}{o m m}{%
    \IfNoValueTF{##1}{\frac{\partial ##2}{\partial ##3}}%
    {\frac{\partial^{##1} ##2}{\partial ##3^{##1}}}%
  }
  \NewDocumentCommand{\dv}{o m m}{%
    \IfNoValueTF{##1}{\frac{\mathrm{d} ##2}{\mathrm{d} ##3}}%
    {\frac{\mathrm{d}^{##1} ##2}{\mathrm{d} ##3^{##1}}}%
  }
  \providecommand{\dd}{\mathop{}\!\mathrm{d}}
  \providecommand{\grad}{\nabla}
  \providecommand{\abs}[1]{\left\lvert ##1\right\rvert}
  \providecommand{\norm}[1]{\left\lVert ##1\right\rVert}
  \providecommand{\ket}[1]{\left\lvert ##1\right\rangle}
  \providecommand{\bra}[1]{\left\langle ##1\right\rvert}
}
\usepackage{booktabs,array,longtable,tabularx}
\usepackage{hyperref}
\usepackage{enumitem}
\usepackage{fancyhdr}
\usepackage{titlesec}
\usepackage{tcolorbox}
\usepackage{ragged2e}
\usepackage{setspace}
\IfFileExists{siunitx.sty}{%
  \usepackage{siunitx}%
}{%
  % Minimal text-safe fallback for the small number of \SI and \si calls in
  % this repository. Full siunitx formatting is used when available.
  \providecommand{\si}[1]{\ensuremath{\mathrm{##1}}}
  \providecommand{\SI}[2]{\ensuremath{##1\,\mathrm{##2}}}
}
\hypersetup{colorlinks=true, linkcolor=blue!50!black, urlcolor=blue!50!black, citecolor=blue!50!black}
\definecolor{TitleBlue}{HTML}{163A5F}
\definecolor{AccentTeal}{HTML}{2D6F73}
\definecolor{SoftBlue}{HTML}{EAF3F8}
\definecolor{SoftTeal}{HTML}{EDF8F6}
\definecolor{SoftGray}{HTML}{F5F7FA}
\definecolor{RuleGray}{HTML}{D7DEE8}
\pagestyle{fancy}
\fancyhf{}
\lhead{RadiationEffects\_proj2}
\rhead{Technical Notes}
\cfoot{\thepage}
\setlength{\headheight}{14pt}
\setlength{\parskip}{0.5em}
\setlength{\parindent}{0pt}
\onehalfspacing
\numberwithin{equation}{section}
\titleformat{\section}{\Large\bfseries\color{TitleBlue}}{\thesection}{0.6em}{}
\titleformat{\subsection}{\large\bfseries\color{AccentTeal}}{\thesubsection}{0.6em}{}
\titleformat{\subsubsection}{\normalsize\bfseries\color{TitleBlue}}{\thesubsubsection}{0.6em}{}
\newtcolorbox{overviewbox}{colback=SoftBlue,colframe=TitleBlue,boxrule=0.8pt,arc=2mm,left=3mm,right=3mm,top=2mm,bottom=2mm}
\newtcolorbox{physicsbox}{colback=SoftTeal,colframe=AccentTeal,boxrule=0.8pt,arc=2mm,left=3mm,right=3mm,top=2mm,bottom=2mm}
\newtcolorbox{notebox}{colback=SoftGray,colframe=AccentTeal,boxrule=0.6pt,arc=2mm,left=3mm,right=3mm,top=2mm,bottom=2mm}
\newcommand{\DocTitle}[3]{%
\begin{center}
{\LARGE \bfseries \color{TitleBlue} #1}\\[0.5em]
{\large \color{AccentTeal} #2}\\[0.8em]
{\normalsize #3}
\end{center}
\vspace{0.5em}}
\newenvironment{symtable}{\rowcolors{2}{SoftGray}{white}\begin{longtable}{>{\RaggedRight\arraybackslash}p{0.18\textwidth} >{\RaggedRight\arraybackslash}p{0.25\textwidth} >{\RaggedRight\arraybackslash}p{0.47\textwidth}}\toprule \textbf{Symbol / acronym} & \textbf{Name} & \textbf{Meaning in this document} \\ \midrule\endfirsthead\toprule \textbf{Symbol / acronym} & \textbf{Name} & \textbf{Meaning in this document} \\ \midrule\endhead}{\bottomrule\end{longtable}}


\newcommand{\ProjectTOC}{%
\vspace{0.6em}
{\hypersetup{linkcolor=TitleBlue,urlcolor=TitleBlue,citecolor=TitleBlue}\tableofcontents}
\vspace{1.0em}}

\newcommand{\SymbolsAcronymsSection}{%
\section*{Symbols and acronyms used in this document}%
\addcontentsline{toc}{section}{Symbols and acronyms used in this document}}

\newcommand{\rvec}{\bm{r}}
\newcommand{\qvec}{\bm{q}}
\newcommand{\nvec}{\bm{n}}
\newcommand{\xvec}{\bm{x}}
\newcommand{\kB}{k_{\mathrm{B}}}
\newcommand{\qp}{\mathrm{qp}}
\newcommand{\JJ}{\mathrm{JJ}}
\newcommand{\mfp}{\Lambda}
\allowdisplaybreaks

\usepackage{listings}
\usepackage{caption}

\rhead{Module 2 FEM Implementation Note}

\definecolor{CodeBackground}{HTML}{F7F9FC}
\definecolor{CodeFrame}{HTML}{C8D3E0}
\definecolor{CodeKeyword}{HTML}{163A5F}
\definecolor{CodeComment}{HTML}{2D6F73}
\definecolor{CodeString}{HTML}{8A4B08}

\lstdefinestyle{projectmatlab}{
  language=Matlab,
  basicstyle=\ttfamily\small,
  keywordstyle=\color{CodeKeyword}\bfseries,
  commentstyle=\color{CodeComment},
  stringstyle=\color{CodeString},
  backgroundcolor=\color{CodeBackground},
  frame=single,
  rulecolor=\color{CodeFrame},
  breaklines=true,
  breakatwhitespace=false,
  columns=fullflexible,
  keepspaces=true,
  showstringspaces=false,
  numbers=left,
  numberstyle=\scriptsize\color{gray},
  numbersep=8pt,
  xleftmargin=1.2em,
  framexleftmargin=1.0em,
  aboveskip=0.8em,
  belowskip=0.8em
}

\lstdefinestyle{projectshell}{
  basicstyle=\ttfamily\small,
  backgroundcolor=\color{CodeBackground},
  frame=single,
  rulecolor=\color{CodeFrame},
  breaklines=true,
  columns=fullflexible,
  keepspaces=true,
  showstringspaces=false,
  xleftmargin=1.2em,
  framexleftmargin=1.0em,
  aboveskip=0.8em,
  belowskip=0.8em
}

\newcommand{\matlabfile}[1]{\path{#1}}
\newcommand{\phivec}{\bm{\Phi}}
\newcommand{\Evec}{\mathbf{E}}

\begin{document}

\DocTitle{Module 2 FEM Implementation Note}{Detailed How-To Guide for \texttt{RadiationEffects\_proj2}}{Running, testing, inspecting, modifying, and diagnosing the two-dimensional electrostatic finite-element solver}

\hrule
\vspace{0.8em}

\begin{overviewbox}
\textbf{Purpose and scope.} This is the operational companion to the Module 2 physics note. It documents the MATLAB finite-element method (FEM) implementation of the two-dimensional electrostatic Poisson problem on both the original rectangular verification domain and the validated Module 9 transmon substrate mesh. It also documents the factorized multi-configuration FEM dataset path beginning at \matlabfile{main_module2_batch_fem_dataset.m}. It explains what to run, how named electrode tags are consumed, which files are called, what each returned field means, how to execute all legacy, transmon, and batch verification tests, and how to diagnose an unexpected result. The separate PINN guide remains authoritative for the neural-network path.
\end{overviewbox}

\begin{notebox}
\textbf{Fastest correct start.} Open MATLAB in the repository's \matlabfile{matlab/} directory, run \texttt{setup\_project\_paths}, then run \texttt{test\_module2\_fem\_all\_2d}. After the suite passes, generate the two geometry-aware demonstrations with \texttt{main\_module2\_electrostatics('transmon\_laplace')} and \texttt{main\_module2\_electrostatics('transmon\_trapped\_charge')}. Then generate the conditional-surrogate pilot data with \texttt{[dataset,report] = main\_module2\_batch\_fem\_dataset()}. The batch MAT-file and validation summary are written under \matlabfile{matlab/outputs/module2_batch_fem_2d/}; plots remain disabled by default.
\end{notebox}

\ProjectTOC

\SymbolsAcronymsSection

\renewcommand{\arraystretch}{1.2}
\setlength{\tabcolsep}{5pt}
\begin{symtable}
FEM & Finite-element method & Galerkin discretization using three-node linear triangular elements. \\
T3 & Three-node triangle & Linear triangular element with one potential degree of freedom at each vertex. \\
$\Omega$ & Computational domain & Either the legacy rectangle $[0,L_x]\times[0,L_y]$ or Module 9 substrate $[-1.5,1.5]\,\mathrm{mm}\times[-0.5,0]\,\mathrm{mm}$. \\
$\phi$ & Electrostatic potential & Primary FEM unknown, stored at mesh nodes in volts. \\
$\Evec$ & Electric field & Postprocessed as $\Evec=-\nabla\phi$, in volts per meter. \\
$\rho$ & Space-charge density & Volumetric source in coulombs per cubic meter. \\
$\varepsilon_{\mathrm{Si}}$ & Silicon permittivity & Product of relative and vacuum permittivity, in farads per meter. \\
$N_a$ & Shape function & Linear basis function for local triangle node $a$. \\
$K$ & Stiffness matrix & Unconstrained sparse global Poisson matrix. \\
$K_{\mathrm{bc}}$ & Boundary-modified matrix & Matrix after strong Dirichlet constraints are imposed. \\
$\mathbf{b}$ & Source vector & Assembled right-hand side before boundary modification. \\
$\phivec$ & Nodal potential vector & Solution of $K_{\mathrm{bc}}\phivec=\mathbf{b}_{\mathrm{bc}}$. \\
$\bm{\mu}$ & Charge-cloud parameters & $[C_{\mathrm{peak}},x_c,z_c,\sigma_x,\sigma_z]$ for one complete batch configuration. \\
BC & Boundary condition & Dirichlet potential or natural homogeneous Neumann field condition. \\
SI & International System of Units & Unit convention used by all Module 2 parameters. \\
\end{symtable}

\section{What the Module 2 FEM code does}

The solver computes the electrostatic potential generated by mobile carriers, ionized dopants, and an effective charged-defect population. The strong-form equation is
\begin{equation}
\nabla\cdot\left(\varepsilon_{\mathrm{Si}}\nabla\phi\right)=-\rho,
\qquad
\Evec=-\nabla\phi.
\end{equation}
The implemented weak form is
\begin{equation}
\int_{\Omega}\varepsilon_{\mathrm{Si}}\nabla w\cdot\nabla\phi\,\dd\Omega
=
\int_{\Omega}w\rho\,\dd\Omega,
\end{equation}
with prescribed potential on selected boundaries and zero normal potential gradient on natural boundaries. The FEM approximation produces
\begin{equation}
K\phivec=\mathbf{b}.
\end{equation}

The current path is a \textbf{steady, linear, single-material electrostatic solve}. There is no time step in Module 2 itself. Time dependence enters a future coupled calculation only because Module 3, Module 4, and Module 5 can supply a new charge state at successive coupled times.

\subsection{Implemented capabilities}

\begin{itemize}[leftmargin=*]
  \item Legacy structured rectangular meshes and exact reuse of the validated nonuniform Module 9 substrate mesh.
  \item Constant silicon permittivity throughout the domain.
  \item Nodal charge density from electrons, holes, ionized dopants, and one effective Gaussian defect population.
  \item A special uniform-charge source used for analytical verification.
  \item Strong constant-valued Dirichlet conditions on rectangular sides or named Module 9 tags.
  \item Natural homogeneous Neumann conditions on sides that are not Dirichlet.
  \item Sparse global assembly, direct sparse solution with MATLAB's backslash operator, element electric fields, and area-weighted nodal electric fields.
  \item Explicit free-node residual, fixed-voltage error, and global source/reaction balance diagnostics.
  \item Transmon field plots with left-electrode, right-electrode, and JJ overlays in $(x,z)$ coordinates.
  \item MAT-file, PNG-figure, and text-summary output through the top-level driver.
  \item Deterministic 17-configuration pilot dataset with one stiffness assembly, one sparse Cholesky factorization, and complete-case train/validation/test splits.
\end{itemize}

\subsection{Important current limitations}

\begin{itemize}[leftmargin=*]
  \item The implementation does \emph{not} assemble a nonzero Neumann boundary integral. The fields \texttt{params.bc.*.dphidn} are descriptive placeholders for the present zero-flux cases. Changing one to a nonzero value does not change the solution.
  \item Permittivity is supplied as one scalar. Spatially varying or tensor permittivity is not yet supported.
  \item Dirichlet values are scalar constants per side or named tag; a spatially varying contact profile is not accepted by the current boundary collector.
  \item The default defect source is generated internally. A nodal defect map from Module 3 is not yet an input option in this standalone source builder.
  \item The solver does not solve carrier statistics or trap occupancy self-consistently. It accepts concentrations and an effective defect charge number as prescribed inputs.
  \item The thin metal films remain zero-thickness boundary tags. This electrostatic solve does not resolve metal or oxide thickness as bulk elements.
  \item The JJ tag is a scoring/interface location, not a resolved Josephson constitutive model.
  \item An all-Neumann problem is singular. The solver now stops explicitly when no Dirichlet nodes are found.
\end{itemize}

\section{Software requirements and starting location}

\subsection{Required software}

The Module 2 FEM path calls core MATLAB operations only: structures, sparse matrices, backslash and sparse Cholesky solves, basic plotting, file I/O, and path management. It does not call the Deep Learning Toolbox and does not require the Module 2 PINN files. Geant4 is also not required for the six standalone FEM cases or the batch dataset generator.

The practical requirements are:
\begin{itemize}[leftmargin=*]
  \item MATLAB with permission to read the repository and write under \matlabfile{matlab/outputs/};
  \item sufficient memory for the requested mesh;
  \item a graphics-capable MATLAB installation when \texttt{makePlots=true}. The figures are created with \texttt{'Visible','off'}, so an interactive figure window is not required in the usual batch configuration.
\end{itemize}

\subsection{Correct working directory}

The least ambiguous starting point is the project's MATLAB directory:
\begin{lstlisting}[style=projectmatlab]
cd('C:/path/to/RadiationEffects_proj2/matlab')
setup_project_paths
\end{lstlisting}
Replace the example path with the actual extracted repository path. Forward slashes work in MATLAB on Windows. If MATLAB was opened directly in \matlabfile{matlab/}, only the second line is needed.

\texttt{setup\_project\_paths.m} adds the following directories:
\begin{itemize}[leftmargin=*]
  \item \matlabfile{matlab/};
  \item \matlabfile{matlab/cases/};
  \item \matlabfile{matlab/tests/};
  \item \matlabfile{matlab/pinn/};
  \item all folders recursively below \matlabfile{matlab/src/}.
\end{itemize}
The PINN directory is added for the complete project, but it is not used by the FEM run documented here.

\subsection{Path preflight check}

Before solving, confirm that MATLAB resolves the intended files:
\begin{lstlisting}[style=projectmatlab]
which main_module2_electrostatics -all
which solve_poisson_defect_space_charge_2d -all
which default_module2_params -all
\end{lstlisting}
Each command should identify the corresponding file inside the current project's \matlabfile{matlab/} tree. Multiple results can indicate that an older project copy is still on the MATLAB path. Remove the stale path or start a fresh MATLAB session before comparing results.

\section{Quick-start execution recipes}

\subsection{Default run}

Calling the driver without an argument selects the localized charged-defect case:
\begin{lstlisting}[style=projectmatlab]
setup_project_paths
result = main_module2_electrostatics;
\end{lstlisting}
This is equivalent to
\begin{lstlisting}[style=projectmatlab]
result = main_module2_electrostatics('localized_defect_charge');
\end{lstlisting}

\subsection{Run all six built-in FEM cases}

\begin{lstlisting}[style=projectmatlab]
setup_project_paths

r0 = main_module2_electrostatics('zero_charge');
r1 = main_module2_electrostatics('linear_potential');
r2 = main_module2_electrostatics('uniform_space_charge');
r3 = main_module2_electrostatics('localized_defect_charge');
r4 = main_module2_electrostatics('transmon_laplace');
r5 = main_module2_electrostatics('transmon_trapped_charge');
\end{lstlisting}
The same output filenames are overwritten when the same case is rerun. Distinct built-in case names do not overwrite one another because the case name prefixes every generated file.

\subsection{Generate the 17-case Module 9 pilot dataset}

The batch generator is a separate user-facing driver. It does not create one result structure or one set of figures per configuration:
\begin{lstlisting}[style=projectmatlab]
setup_project_paths
test_module2_batch_fem_dataset_2d
[dataset, report] = main_module2_batch_fem_dataset();
\end{lstlisting}
The default call saves one compact MAT-file and one text manifest. It creates no PNGs and opens no figures. To export only four explicitly configured audit cases, opt in deliberately:
\begin{lstlisting}[style=projectmatlab]
options = default_module2_batch_fem_options_2d();
options.makeRepresentativePlots = true;
[dataset, report] = main_module2_batch_fem_dataset(options);
\end{lstlisting}

\subsection{Run one case from a system terminal}

MATLAB's batch interface can run the solver without first opening the desktop:
\begin{lstlisting}[style=projectshell]
matlab -batch "cd('C:/path/to/RadiationEffects_proj2/matlab'); setup_project_paths; result = main_module2_electrostatics('localized_defect_charge');"
\end{lstlisting}
The terminal process returns a nonzero exit status if MATLAB encounters an uncaught error, which is useful for automated verification.

\subsection{Run without writing plots or a MAT-file}

The top-level driver constructs its own default parameter structure, so output flags cannot be passed through its function signature. For a diagnostic solve with reduced file output, call the parameter factory and solver directly:
\begin{lstlisting}[style=projectmatlab]
setup_project_paths
params = default_module2_params('localized_defect_charge');
params.makePlots = false;
params.saveMat = false;
result = solve_poisson_defect_space_charge_2d(params);
\end{lstlisting}
In this direct-solver path, \texttt{makePlots} and \texttt{saveMat} are not read because plotting and saving belong to the top-level driver. The flags are still useful when the same \texttt{params} structure is later passed through a custom driver. The direct call itself writes no files unless plotting and summary functions are invoked explicitly.

\section{Built-in cases and expected behavior}

\begin{longtable}{p{0.24\textwidth} p{0.28\textwidth} p{0.40\textwidth}}
\toprule
\textbf{Case name} & \textbf{Configuration} & \textbf{Expected result and purpose} \\
\midrule
\endfirsthead
\toprule
\textbf{Case name} & \textbf{Configuration} & \textbf{Expected result and purpose} \\
\midrule
\endhead
\matlabfile{zero_charge} & $\rho=0$; left and right potentials both zero. & $\phi=0$ and $\Evec=0$ to numerical precision. This detects unintended sources or boundary values. \\
\matlabfile{linear_potential} & $\rho=0$; left potential $0$ V and right potential $1$ V. & Exact Laplace solution $\phi=x/L_x$ with uniform $E_x=-1/L_x$. This checks mesh assembly, Dirichlet insertion, and field sign. \\
\matlabfile{uniform_space_charge} & Explicit $\rho_{0}=5.0\times10^{-4}$ C/m$^3$; left and right potentials zero. & Nearly one-dimensional quadratic solution $\phi=\rho_0x(L_x-x)/(2\varepsilon_{\mathrm{Si}})$. This checks the source vector. \\
\matlabfile{localized_defect_charge} & Positive Gaussian effective defect population centered in the device; grounded left and right contacts. & Finite positive potential perturbation and localized positive charge. This is a physical sanity case, not an exact-solution verification. \\
\matlabfile{transmon_laplace} & Module 9 mesh; $\rho=0$; $0$ V on \texttt{left\_electrode} and $1$ mV on \texttt{right\_electrode}; all other boundaries natural. & Mixed-boundary Laplace field on the chip cross-section. This verifies named-tag constraints, electrode separation, maximum range, residual, and reaction balance. \\
\matlabfile{transmon_trapped_charge} & Module 9 mesh; both named electrodes grounded; positive Gaussian defect charge centered $20\,\mu$m below the JJ. & Positive localized potential perturbation on the shared transmon mesh. This tests charge mapping, tagged grounding, residual, and global source/reaction balance. \\
\bottomrule
\end{longtable}

The files in \matlabfile{matlab/cases/} are small convenience functions that return these parameter structures. The main driver does not call the case functions; it calls \texttt{default\_module2\_params(caseName)} directly. For example,
\begin{lstlisting}[style=projectmatlab]
params = case_module2_linear_potential_2d();
result = solve_poisson_defect_space_charge_2d(params);
\end{lstlisting}
is a valid direct-solver route, but it bypasses the driver's automatic saving, plotting, summary writing, and console report.

\section{File map and responsibility of each component}

\begin{longtable}{p{0.39\textwidth} p{0.53\textwidth}}
\toprule
\textbf{File} & \textbf{Responsibility} \\
\midrule
\endfirsthead
\toprule
\textbf{File} & \textbf{Responsibility} \\
\midrule
\endhead
\matlabfile{matlab/main_module2_electrostatics.m} & User-facing driver. Selects a built-in case, runs the solver, creates the output directory, saves results, makes plots, writes a summary, and prints scalar diagnostics. \\
\matlabfile{matlab/main_module2_batch_fem_dataset.m} & User-facing multi-configuration driver. Generates, validates, and saves one Module 9-aware dataset; representative plots are explicit opt-in outputs. \\
\matlabfile{matlab/setup_project_paths.m} & Adds the driver, cases, tests, PINN directory, and all source subdirectories to the MATLAB path. \\
\matlabfile{matlab/src/module2_electrostatics/default_module2_params.m} & Defines constants, geometry mode, source parameters, side/tag boundary conditions, output controls, and six supported case modifications. \\
\matlabfile{matlab/src/module2_electrostatics/solve_poisson_defect_space_charge_2d.m} & Numerical orchestrator. Calls mesh generation, charge evaluation, assembly, boundary processing, sparse solve, and field postprocessing. \\
\matlabfile{matlab/src/module2_electrostatics/resolve_module2_mesh_2d.m} & Selects the legacy rectangle or validates and reuses a supplied Module 9 geometry; synchronizes domain and substrate permittivity metadata. \\
\matlabfile{matlab/src/module2_electrostatics/make_rectangular_tri_mesh_2d.m} & Creates nodes, counter-clockwise T3 connectivity, coordinate vectors, dimensions, and left/right/top/bottom boundary-node sets. \\
\matlabfile{matlab/src/module2_electrostatics/build_space_charge_module2_2d.m} & Evaluates the nodal source $\rho$. Uses explicit \texttt{rho\_uniform} when present; otherwise evaluates carriers, dopants, and a Gaussian effective defect population. \\
\matlabfile{matlab/src/module2_electrostatics/assemble_poisson_fem_2d.m} & Loops over triangles, forms local stiffness and source contributions, scatters them into a sparse global matrix and right-hand side, and stores element areas and centroids. \\
\matlabfile{matlab/src/module2_electrostatics/assemble_poisson_source_rhs_2d.m} & Reuses stored triangle areas to assemble only the new source vector for each batch configuration. \\
\matlabfile{matlab/src/module2_electrostatics/compute_element_stiffness_triangle.m} & Calculates signed area, shape-function gradients, and $K_e=\varepsilon A_e(\nabla N)(\nabla N)^T$. Rejects non-positive element orientation. \\
\matlabfile{matlab/src/module2_electrostatics/compute_element_source_triangle.m} & Calculates the consistent T3 source vector using the triangle mass-like matrix multiplied by the three nodal charge values. \\
\matlabfile{matlab/src/module2_electrostatics/get_module2_dirichlet_nodes.m} & Collects nodes and scalar values from rectangle sides and named Module 9 tags. Rejects missing, empty, or mutually conflicting constraints. \\
\matlabfile{matlab/src/module2_electrostatics/apply_dirichlet_bc.m} & Adjusts the right-hand side for known values, zeros constrained rows and columns, and inserts unit diagonal equations. \\
\matlabfile{matlab/src/module2_electrostatics/compute_electric_field_from_potential_2d.m} & Computes elementwise $-\nabla\phi$ and area-weighted nodal averages. \\
\matlabfile{matlab/src/module2_electrostatics/compute_module2_fem_diagnostics_2d.m} & Computes free-node residual, Dirichlet error, constraint reactions, and discrete global source/reaction balance. \\
\matlabfile{matlab/src/module2_electrostatics/plot_module2_result_2d.m} & Writes potential, space-charge, and electric-field-magnitude PNG figures. \\
\matlabfile{matlab/src/module2_electrostatics/write_module2_summary.m} & Writes geometry, mesh, material, source, potential, field, and boundary-node scalar values to a text file. \\
\matlabfile{matlab/src/module2_electrostatics/default_module2_batch_fem_options_2d.m} & Defines dataset name, design, fixed electrode voltages, output controls, representative cases, and validation tolerances. \\
\matlabfile{matlab/src/module2_electrostatics/make_module2_batch_parameter_design_2d.m} & Defines the deterministic 17-case Gaussian parameter matrix and complete-case split indices. \\
\matlabfile{matlab/src/module2_electrostatics/generate_module2_batch_fem_dataset_2d.m} & Builds one operator, performs one sparse Cholesky factorization, solves all source vectors, and returns the compact dataset contract without file I/O or plotting. \\
\matlabfile{matlab/src/module2_electrostatics/validate_module2_batch_fem_dataset_2d.m} & Checks schema, shapes, finiteness, disjoint complete-case splits, tagged voltages, JJ treatment, residuals, balance, and factorization reuse. \\
\matlabfile{matlab/src/module2_electrostatics/write_module2_batch_fem_summary.m} & Writes the dataset manifest, split membership, solver counts, and worst diagnostics. \\
\matlabfile{matlab/src/module2_electrostatics/plot_module2_batch_fem_case_2d.m} & Saves one selected charge/potential audit figure; never called by the numerical generator. \\
\matlabfile{matlab/cases/case_module2_*.m} & Convenience parameter factories for all six built-in cases. \\
\matlabfile{matlab/tests/test_module2_*.m} & Standalone FEM verification and sanity tests. PINN and CPINN tests with similar prefixes are outside this note's scope. \\
\bottomrule
\end{longtable}

\section{Exact execution sequence}

For \texttt{main\_module2\_electrostatics('localized\_defect\_charge')}, the call sequence is:
\begin{enumerate}[leftmargin=*]
  \item The driver checks the input and supplies \texttt{'localized\_defect\_charge'} when no name is given.
  \item \texttt{default\_module2\_params} creates the complete parameter structure and applies the selected case branch.
  \item \texttt{resolve\_module2\_mesh\_2d} creates the legacy rectangle or validates/reuses the Module 9 geometry.
  \item \texttt{build\_space\_charge\_module2\_2d} evaluates one charge value per node.
  \item \texttt{assemble\_poisson\_fem\_2d} forms unconstrained \texttt{K} and \texttt{rhs}.
  \item \texttt{get\_module2\_dirichlet\_nodes} resolves side- or tag-based constrained nodes and rejects conflicting values.
  \item \texttt{apply\_dirichlet\_bc} forms \texttt{Kbc} and \texttt{rhsbc}.
  \item MATLAB solves \texttt{phi = Kbc \textbackslash{} rhsbc}.
  \item \texttt{compute\_electric\_field\_from\_potential\_2d} computes element and nodal fields.
  \item \texttt{compute\_module2\_fem\_diagnostics\_2d} evaluates the unconstrained free-node residual, exact boundary error, and global reaction balance.
  \item The solver returns a \texttt{result} structure containing inputs, selected geometry, matrices, solution, fields, boundary provenance, and diagnostics.
  \item The driver saves the MAT-file, creates three figures, writes the text summary, and prints a completion report.
\end{enumerate}

This separation is useful while debugging: call the solver directly to isolate numerical behavior, and call the driver when testing the complete input/output workflow.

\section{Batch FEM dataset contract and execution sequence}

The pilot parameterization is
\begin{equation}
\bm{\mu}=\left[C_{\mathrm{peak}},x_c,z_c,\sigma_x,\sigma_z\right],
\end{equation}
for one elliptical Gaussian effective-defect cloud. The deterministic design contains a zero control and 16 charged cases spanning amplitude, depth, lateral displacement, and anisotropic width. Its 11 training, three validation, and three test entries are indices of \emph{complete configurations}; nodes from a held-out configuration are never placed in the training set.

The efficient path is:
\begin{enumerate}[leftmargin=*]
  \item build and validate one Module 9 geometry;
  \item resolve the left and right electrode nodes once, with both voltages grounded for this first dataset;
  \item assemble the silicon stiffness matrix $K$ once and apply the tagged boundary rows and columns once;
  \item compute one sparse Cholesky factor $K_{\mathrm{bc}}=LL^T$;
  \item for each row of $\bm{\mu}$, evaluate nodal $\rho$, assemble only the source vector, apply the fixed-value right-hand-side correction, and solve $L\mathbf{y}=\mathbf{b}_{\mathrm{bc}}$ followed by $L^T\phivec=\mathbf{y}$;
  \item record compact scalar residual, boundary, balance, and range diagnostics;
  \item validate the aggregate contract and save one MAT-file.
\end{enumerate}

The saved structure stores \texttt{mesh} and \texttt{tags} once. The matrices \texttt{rho} and \texttt{phi} contain one row per mesh node and one column per complete case. The matrix \texttt{parameters} contains one row per case and five columns; its column names and units are stored separately. The fields \texttt{trainCases}, \texttt{validationCases}, and \texttt{testCases} identify whole columns. The stiffness matrix, Cholesky factor, and full residual vectors are intentionally omitted because the conditional PINN does not need them and they would inflate every saved dataset.

\begin{notebox}
\textbf{Linearity decision.} The pilot keeps both electrode voltages fixed at zero. For the current linear Poisson equation, nonzero contact-voltage contributions can later be supplied exactly through precomputed Laplace basis solutions and superposition. The first conditional PINN therefore learns the defect-charge response without wasting capacity on a contribution already known analytically.
\end{notebox}

\section{Parameter reference}

All physical inputs are in SI units unless noted otherwise. The default values below are defined before the case-specific switch is evaluated.

\subsection{Physical constants and material}

\begin{longtable}{p{0.25\textwidth} p{0.18\textwidth} p{0.18\textwidth} p{0.30\textwidth}}
\toprule
\textbf{Field} & \textbf{Default} & \textbf{Units} & \textbf{Meaning} \\
\midrule
\endfirsthead
\toprule
\textbf{Field} & \textbf{Default} & \textbf{Units} & \textbf{Meaning} \\
\midrule
\endhead
\texttt{q} & $1.602176634\times10^{-19}$ & C & Elementary charge magnitude. \\
\texttt{eps0} & $8.8541878128\times10^{-12}$ & F/m & Vacuum permittivity. \\
\texttt{eps\_rel\_si} & $11.7$ & dimensionless & Relative permittivity of silicon. \\
\texttt{eps\_si} & \matlabfile{eps_rel_si*eps0} & F/m & Scalar permittivity passed to element assembly. Recompute it after changing either factor. \\
\bottomrule
\end{longtable}

\subsection{Geometry and mesh}

\begin{longtable}{p{0.25\textwidth} p{0.18\textwidth} p{0.18\textwidth} p{0.30\textwidth}}
\toprule
\textbf{Field} & \textbf{Default} & \textbf{Units} & \textbf{Meaning} \\
\midrule
\endfirsthead
\toprule
\textbf{Field} & \textbf{Default} & \textbf{Units} & \textbf{Meaning} \\
\midrule
\endhead
\texttt{Lx} & $10\times10^{-6}$ & m & Domain length in $x$. \\
\texttt{Ly} & $4\times10^{-6}$ & m & Domain length in $y$. \\
\texttt{nx} & $41$ & nodes & Number of structured node locations along $x$. Must be at least two. \\
\texttt{ny} & $21$ & nodes & Number of structured node locations along $y$. Must be at least two. \\
\texttt{geometryMode} & \texttt{rectangle} & -- & Selects \texttt{rectangle} or \texttt{module9\_transmon}. The transmon cases set this automatically. \\
\texttt{module9Geometry} & empty & structure & Optional prebuilt geometry. When supplied, Module 2 reuses its node coordinates, connectivity, and tags exactly. \\
\texttt{coordinateNames} & \texttt{\{x,y\}} & -- & Becomes \texttt{\{x,z\}} for Module 9 cases. \\
\bottomrule
\end{longtable}
The node count is $N= n_xn_y$, and the triangle count is
\begin{equation}
N_e=2(n_x-1)(n_y-1).
\end{equation}
Increasing \texttt{nx} and \texttt{ny} improves spatial resolution but increases matrix size, assembly time, solution time, and output size.

For \texttt{module9\_transmon}, \texttt{nx} and \texttt{ny} are synchronized from the accepted nonuniform geometry and the legacy rectangular mesh generator is bypassed. To prove exact shared-mesh reuse:
\begin{lstlisting}[style=projectmatlab]
shared = build_module9_transmon_geometry_2d();
p = default_module2_params('transmon_laplace');
p.module9Geometry = shared;
r = solve_poisson_defect_space_charge_2d(p);
assert(isequal(r.mesh.nodes, shared.mesh.nodes))
assert(isequal(r.mesh.elems, shared.mesh.elems))
\end{lstlisting}

\subsection{Carrier, dopant, and defect charge}

\begin{longtable}{p{0.27\textwidth} p{0.18\textwidth} p{0.16\textwidth} p{0.30\textwidth}}
\toprule
\textbf{Field} & \textbf{Default} & \textbf{Units} & \textbf{Meaning} \\
\midrule
\endfirsthead
\toprule
\textbf{Field} & \textbf{Default} & \textbf{Units} & \textbf{Meaning} \\
\midrule
\endhead
\texttt{n} & $0$ & m$^{-3}$ & Uniform electron concentration, entering with a negative sign. \\
\texttt{p} & $0$ & m$^{-3}$ & Uniform hole concentration, entering with a positive sign. \\
\texttt{ND\_plus} & $0$ & m$^{-3}$ & Uniform ionized donor concentration. \\
\texttt{NA\_minus} & $0$ & m$^{-3}$ & Uniform ionized acceptor concentration. \\
\texttt{Cdef\_background} & $0$ & m$^{-3}$ & Uniform effective defect concentration. \\
\texttt{Cdef\_peak} & $2.0\times10^{20}$ & m$^{-3}$ & Peak Gaussian defect concentration before the case switch. \\
\texttt{Cdef\_x0} & $0.5L_x$ & m & Gaussian center in $x$. \\
\texttt{Cdef\_y0} & $0.5L_y$ & m & Gaussian center in $y$. \\
\texttt{Cdef\_sigma\_x} & $0.10L_x$ & m & Gaussian standard deviation in $x$. Must be positive. \\
\texttt{Cdef\_sigma\_y} & $0.15L_y$ & m & Gaussian standard deviation in $y$. Must be positive. \\
\texttt{zdef} & $1.0$ & dimensionless & Effective charge number per defect. Positive gives positive defect charge. \\
\texttt{rho\_uniform} & case-specific & C/m$^3$ & When this field exists, it bypasses all carrier, dopant, and defect calculations and supplies an explicit uniform source. \\
\bottomrule
\end{longtable}

Unless \texttt{rho\_uniform} exists, the source builder evaluates
\begin{align}
C_{\mathrm{def}}(x,y)
&=C_{\mathrm{bg}}+C_{\mathrm{peak}}
\exp\left[-\frac{1}{2}\left(\frac{x-x_0}{\sigma_x}\right)^2
-\frac{1}{2}\left(\frac{y-y_0}{\sigma_y}\right)^2\right],\\
\rho(x,y)
&=q\left[p-n+N_D^+-N_A^-+z_{\mathrm{def}}C_{\mathrm{def}}(x,y)\right].
\end{align}

\begin{notebox}
\textbf{Common source pitfall.} If a parameter structure was created for \texttt{uniform\_space\_charge}, deleting or changing carrier/defect fields will have no effect while \texttt{rho\_uniform} still exists. Remove it with \texttt{params = rmfield(params,'rho\_uniform')} before using the general source expression.
\end{notebox}

\subsection{Boundary conditions}

Each side has a \texttt{type}. A Dirichlet side also has \texttt{value}; a nominal Neumann side has \texttt{dphidn}. The defaults are grounded left and right contacts with natural top and bottom boundaries:
\begin{lstlisting}[style=projectmatlab]
params.bc.left.type = 'dirichlet';
params.bc.left.value = 0.0;
params.bc.right.type = 'dirichlet';
params.bc.right.value = 0.0;
params.bc.bottom.type = 'neumann';
params.bc.bottom.dphidn = 0.0;
params.bc.top.type = 'neumann';
params.bc.top.dphidn = 0.0;
\end{lstlisting}
Only the \texttt{'dirichlet'} branches are explicitly processed. All other sides contribute the natural homogeneous boundary term of the weak form. A nonzero \texttt{dphidn} value is not yet integrated.

For transmon cases, \matlabfile{params.bc.named} points to \matlabfile{electrode.left_electrode} and \matlabfile{electrode.right_electrode}. The collector combines side and named-tag nodes, removes compatible duplicates, and throws an error if the same node receives incompatible voltages. It never applies a boundary condition to \matlabfile{jj_sensitive}.

\subsection{Output controls}

\begin{longtable}{p{0.26\textwidth} p{0.18\textwidth} p{0.48\textwidth}}
\toprule
\textbf{Field} & \textbf{Default} & \textbf{Effect in the top-level driver} \\
\midrule
\endfirsthead
\toprule
\textbf{Field} & \textbf{Default} & \textbf{Effect in the top-level driver} \\
\midrule
\endhead
\texttt{outputDir} & \matlabfile{outputs/module2_2d} & Resolved relative to the directory containing the top-level driver. \\
\texttt{makePlots} & \texttt{true} & Calls \texttt{plot\_module2\_result\_2d}. \\
\texttt{saveMat} & \texttt{true} & Saves the complete \texttt{result} structure. \\
\bottomrule
\end{longtable}
The text summary is always written by the top-level driver, independent of \texttt{makePlots} and \texttt{saveMat}.

\section{Custom runs without editing the baseline defaults}

For exploratory work, create a parameter structure, override fields in the command window or a new script, and call the numerical solver directly. This preserves the validated defaults.

\subsection{Mesh-refined localized source}

\begin{lstlisting}[style=projectmatlab]
setup_project_paths
params = default_module2_params('localized_defect_charge');
params.nx = 81;
params.ny = 41;
params.Cdef_peak = 5.0e20;
params.zdef = -1.0;
result = solve_poisson_defect_space_charge_2d(params);

fprintf('Potential range: [%g, %g] V\n', ...
    min(result.phi), max(result.phi));
\end{lstlisting}
Changing \texttt{zdef} from $+1$ to $-1$ reverses the defect-charge contribution. With symmetric grounded contacts and no other charge, the potential perturbation should reverse sign.

\subsection{Uniform dopant-only source}

\begin{lstlisting}[style=projectmatlab]
params = default_module2_params('zero_charge');
params.ND_plus = 1.0e20;
params.NA_minus = 0.0;
result = solve_poisson_defect_space_charge_2d(params);
\end{lstlisting}
Because the concentrations are treated as prescribed uniform values, this is a source-driven electrostatic experiment, not a self-consistent semiconductor equilibrium calculation.

\subsection{Saving a direct-solver result deliberately}

\begin{lstlisting}[style=projectmatlab]
outputDir = fullfile(pwd, 'outputs', 'module2_2d', 'custom_probe');
if ~exist(outputDir, 'dir')
    mkdir(outputDir);
end
save(fullfile(outputDir, 'custom_probe_results.mat'), 'result');
plot_module2_result_2d(result, outputDir);
write_module2_summary(result, outputDir);
\end{lstlisting}
Assign a unique \texttt{params.caseName} before plotting or summary writing if several custom cases share the same output directory.

\section{Returned data structure: what to inspect}

The numerical solver returns a structure with the following top-level fields.

\begin{longtable}{p{0.27\textwidth} p{0.65\textwidth}}
\toprule
\textbf{Field} & \textbf{Contents} \\
\midrule
\endfirsthead
\toprule
\textbf{Field} & \textbf{Contents} \\
\midrule
\endhead
\texttt{result.params} & Complete parameter structure used for the solve. \\
\texttt{result.mesh} & Nodes, element connectivity, coordinate vectors, dimensions, domain lengths, and four boundary-node sets. \\
\texttt{result.geometry} & Validated Module 9 geometry, tags, materials, and schema for transmon cases; empty for legacy rectangles. \\
\texttt{result.rho} & $N\times1$ nodal charge-density vector. \\
\texttt{result.K} & Unconstrained sparse stiffness matrix. \\
\texttt{result.rhs} & Unconstrained assembled source vector. \\
\texttt{result.Kbc} & Sparse matrix after Dirichlet insertion. \\
\texttt{result.rhsbc} & Boundary-modified right-hand side. \\
\texttt{result.fixedNodes} & Constrained global node indices. \\
\texttt{result.fixedValues} & Prescribed potentials corresponding to \texttt{fixedNodes}. \\
\texttt{result.bcInfo} & Provenance for each side- or named-tag constraint and the final fixed-node count. \\
\texttt{result.phi} & Solved nodal potential vector. \\
\texttt{result.field} & Element and nodal electric-field arrays, magnitudes, centroids, and areas. \\
\texttt{result.elementData} & Element areas and centroids recorded during assembly. \\
\texttt{result.diagnostics} & Free-node residual, Dirichlet error, constraint reactions, assembled source, and global balance. \\
\texttt{result.maxAbsPhi} & Maximum absolute nodal potential. \\
\texttt{result.maxAbsE} & Maximum area-averaged nodal field magnitude. \\
\bottomrule
\end{longtable}

Useful mesh subfields include \texttt{nodes} ($N\times2$ coordinates), \texttt{elems} ($N_e\times3$ connectivity), \texttt{x}, \texttt{y}, \texttt{nx}, \texttt{ny}, and \texttt{boundary.left/right/top/bottom}. Useful field subfields include \texttt{Eelem}, \texttt{Ex\_elem}, \texttt{Ey\_elem}, \texttt{Enodal}, \texttt{Ex\_nodal}, \texttt{Ey\_nodal}, \texttt{Emag\_elem}, and \texttt{Emag\_nodal}.

\section{How to probe a run numerically}

\subsection{Dimensions, ranges, and finiteness}

\begin{lstlisting}[style=projectmatlab]
size(result.mesh.nodes)
size(result.mesh.elems)
size(result.Kbc)

[min(result.rho), max(result.rho)]
[min(result.phi), max(result.phi)]
[min(result.field.Ex_nodal), max(result.field.Ex_nodal)]
[min(result.field.Ey_nodal), max(result.field.Ey_nodal)]

assert(all(isfinite(result.rho)))
assert(all(isfinite(result.phi)))
assert(all(isfinite(result.field.Enodal), 'all'))
\end{lstlisting}
For MATLAB releases that do not accept \texttt{'all'} in this reduction, replace the final assertion with the equivalent vectorized form:
\begin{lstlisting}[style=projectmatlab]
assert(all(isfinite(result.field.Enodal(:))))
\end{lstlisting}

\subsection{Linear-system residual}

After boundary insertion, the computed solution should satisfy the modified system to floating-point accuracy:
\begin{lstlisting}[style=projectmatlab]
algebraicResidual = result.Kbc * result.phi - result.rhsbc;
relativeResidual = norm(algebraicResidual, 2) / ...
    max(norm(result.rhsbc, 2), eps);
fprintf('Relative algebraic residual = %.3e\n', relativeResidual);
\end{lstlisting}
This checks whether the linear system was solved. It does \emph{not} by itself prove that the PDE, source, or boundary conditions were assembled correctly.

\subsection{Boundary-condition error}

\begin{lstlisting}[style=projectmatlab]
bcError = max(abs(result.phi(result.fixedNodes) - ...
                  result.fixedValues));
fprintf('Maximum Dirichlet error = %.3e V\n', bcError);
\end{lstlisting}

The solver computes this quantity and the two conservation diagnostics automatically:
\begin{lstlisting}[style=projectmatlab]
result.diagnostics.freeResidualRelative
result.diagnostics.dirichletErrorInf
result.diagnostics.globalBalanceRelative
result.bcInfo.entries
\end{lstlisting}
The global check sums the Dirichlet constraint reactions and the assembled volume source. With homogeneous natural boundaries, that sum should be zero to solver tolerance.

\subsection{Matrix structure and conditioning}

\begin{lstlisting}[style=projectmatlab]
figure; spy(result.Kbc); title('Boundary-modified FEM sparsity');
symmetryDefect = norm(result.Kbc - result.Kbc.', 'fro') / ...
    max(norm(result.Kbc, 'fro'), eps);
conditionEstimate = condest(result.Kbc);
\end{lstlisting}
The unconstrained electrostatic stiffness matrix is symmetric. The boundary routine zeros both constrained rows and columns, so \texttt{Kbc} should also be symmetric to roundoff. A very large condition estimate can accompany poor scaling, extreme geometry aspect ratio, or excessive mesh refinement.

\subsection{Probe the node nearest a physical coordinate}

\begin{lstlisting}[style=projectmatlab]
target = [0.50*result.params.Lx, 0.50*result.params.Ly];
distance2 = sum((result.mesh.nodes - target).^2, 2);
[~, nodeID] = min(distance2);

fprintf('Node %d at (%g,%g) m\n', nodeID, ...
    result.mesh.nodes(nodeID,1), result.mesh.nodes(nodeID,2));
fprintf('rho = %g C/m^3, phi = %g V\n', ...
    result.rho(nodeID), result.phi(nodeID));
fprintf('E = (%g,%g) V/m\n', ...
    result.field.Ex_nodal(nodeID), ...
    result.field.Ey_nodal(nodeID));
\end{lstlisting}

\subsection{Inspect one triangle}

\begin{lstlisting}[style=projectmatlab]
elementID = 1;
localNodes = result.mesh.elems(elementID,:);
coordinates = result.mesh.nodes(localNodes,:);
[Ke, area, gradN] = compute_element_stiffness_triangle( ...
    coordinates, result.params.eps_si);
localRho = result.rho(localNodes);
fe = compute_element_source_triangle(area, localRho);
\end{lstlisting}
This is the most direct way to study the connection between triangle geometry, basis gradients, element matrix, source values, and global assembly.

\section{Generated outputs and how to check them}

For a case named \texttt{<case>}, the default driver writes:
\begin{itemize}[leftmargin=*]
  \item \matlabfile{<case>_results.mat}: complete \texttt{result} structure;
  \item \matlabfile{<case>_potential.png}: nodal potential on the triangular mesh;
  \item \matlabfile{<case>_space_charge.png}: nodal space-charge density;
  \item \matlabfile{<case>_electric_field_magnitude.png}: area-averaged nodal field magnitude;
  \item \matlabfile{<case>_summary.txt}: scalar run summary.
\end{itemize}

Legacy outputs use \matlabfile{matlab/outputs/module2_2d/}; the two transmon cases use \matlabfile{matlab/outputs/module2_transmon_2d/}. Transmon figures use millimetres, label the vertical coordinate $z$, and overlay the two electrode groups plus the JJ-sensitive interval. After a run, check:
\begin{enumerate}[leftmargin=*]
  \item the summary reports the requested domain and mesh;
  \item the number of nodes is \texttt{nx*ny};
  \item the number of triangles is \texttt{2*(nx-1)*(ny-1)};
  \item charge, potential, and field ranges are finite;
  \item the potential agrees with the prescribed left and right values;
  \item plot symmetry or sign agrees with the selected source;
  \item the MAT-file loads and reproduces the scalar values in the summary.
\end{enumerate}

\begin{lstlisting}[style=projectmatlab]
loaded = load(fullfile('outputs', 'module2_2d', ...
    'localized_defect_charge_results.mat'));
loaded.result.maxAbsE
\end{lstlisting}

The batch driver uses \matlabfile{matlab/outputs/module2_batch_fem_2d/} and writes:
\begin{itemize}[leftmargin=*]
  \item \matlabfile{module2_transmon_gaussian_batch_v1.mat}: complete dataset structure;
  \item \matlabfile{module2_transmon_gaussian_batch_summary.txt}: case manifest, split membership, operator reuse counts, and worst diagnostics;
  \item \matlabfile{<case>_batch_fem_audit.png}: only when representative plotting is explicitly enabled.
\end{itemize}
Load and inspect the aggregate contract with:
\begin{lstlisting}[style=projectmatlab]
loaded = load(fullfile('outputs', 'module2_batch_fem_2d', ...
    'module2_transmon_gaussian_batch_v1.mat'));
size(loaded.dataset.parameters)   % expected: 17 by 5
size(loaded.dataset.phi)          % expected: nNodes by 17
loaded.dataset.trainCases
loaded.dataset.validationCases
loaded.dataset.testCases
\end{lstlisting}

\section{Verification tests}

\subsection{Run the FEM-only test suite}

From \matlabfile{matlab/}:
\begin{lstlisting}[style=projectmatlab]
setup_project_paths
test_module2_fem_all_2d
\end{lstlisting}
This master runner executes the four legacy tests, the four focused Module 2-on-Module 9 integration tests, and the five-component batch-dataset suite. For a focused diagnosis, use:
\begin{itemize}[leftmargin=*]
  \item \matlabfile{test_module2_transmon_integration_2d} for the geometry-aware single-case path;
  \item \matlabfile{test_module2_batch_fem_dataset_2d} for factorized dataset generation.
\end{itemize}
The following similarly named tests validate different code paths and are outside this FEM test group:
\begin{itemize}[leftmargin=*]
  \item \matlabfile{test_module2_pinn_electrostatics};
  \item \matlabfile{test_module2_cpinn_electrostatics}.
\end{itemize}
Those tests can require the Deep Learning Toolbox.

To stop on the first failure while running the FEM group programmatically:
\begin{lstlisting}[style=projectmatlab]
femTests = { ...
    @test_module2_zero_charge_2d, ...
    @test_module2_linear_potential_2d, ...
    @test_module2_uniform_space_charge_2d, ...
    @test_module2_localized_defect_charge_2d, ...
    @test_module2_transmon_tagged_bc_2d, ...
    @test_module2_transmon_trapped_charge_2d, ...
    @test_module2_transmon_mesh_reuse_2d, ...
    @test_module2_transmon_boundary_guards_2d, ...
    @test_module2_batch_fem_dataset_2d};

for k = 1:numel(femTests)
    feval(femTests{k});
end
\end{lstlisting}

\subsection{What each test proves}

\begin{longtable}{p{0.34\textwidth} p{0.22\textwidth} p{0.35\textwidth}}
\toprule
\textbf{Test} & \textbf{Pass criterion} & \textbf{Interpretation} \\
\midrule
\endfirsthead
\toprule
\textbf{Test} & \textbf{Pass criterion} & \textbf{Interpretation} \\
\midrule
\endhead
\matlabfile{test_module2_zero_charge_2d} & $\max|\phi|<10^{-12}$ V & Zero source and zero contact potential produce the zero solution. \\
\matlabfile{test_module2_linear_potential_2d} & Maximum nodal error below $10^{-10}$ V & Linear T3 basis reproduces the exact charge-free linear solution. \\
\matlabfile{test_module2_uniform_space_charge_2d} & Relative maximum potential error below $5\times10^{-3}$ & Source assembly approaches the analytical one-dimensional quadratic solution on the default mesh. \\
\matlabfile{test_module2_localized_defect_charge_2d} & Finite potential, positive maximum source, and positive maximum potential & Localized positive charge produces a physically plausible response. This is a sanity test, not a convergence proof. \\
\matlabfile{test_module2_transmon_tagged_bc_2d} & Exact $0/1$ mV tag voltages, disjoint electrodes, bounded potential, residual and balance below $10^{-9}$, and three PNG smoke outputs & Module 2 consumes the named electrode tags without shorting through the JJ and all geometry-aware plot paths execute. \\
\matlabfile{test_module2_transmon_trapped_charge_2d} & Positive finite response, exact grounded tags, residual and balance below $10^{-9}$ & Localized defect charge is evaluated on the Module 9 mesh and drives the expected sign. \\
\matlabfile{test_module2_transmon_mesh_reuse_2d} & Node coordinates, connectivity, and electrode edge IDs match exactly & A prebuilt shared Module 9 object is passed through without remeshing or retagging. \\
\matlabfile{test_module2_transmon_boundary_guards_2d} & Unknown tag and conflicting-voltage error identifiers are caught & Misspelled tags cannot fall back silently, and one node cannot receive incompatible constraints. \\
\matlabfile{test_module2_batch_dataset_contract_2d} & All aggregate validation checks pass & Dataset shapes, fields, tags, voltages, scalar diagnostics, and public $(x,z)$ parameter naming satisfy the saved contract. \\
\matlabfile{test_module2_batch_factorization_reuse_2d} & One stiffness assembly, one factorization, and one solve per case & The batch path reuses the operator rather than repeating the expensive setup work. \\
\matlabfile{test_module2_batch_single_case_equivalence_2d} & Potential agrees with the validated single-case path to relative tolerance $10^{-11}$ & Factorized reuse changes execution strategy without changing the electrostatic solution. \\
\matlabfile{test_module2_batch_split_integrity_2d} & Splits are disjoint and cover every case once & No configuration-level leakage is introduced before conditional-PINN training. \\
\matlabfile{test_module2_batch_reproducible_design_2d} & Repeated design construction is identical & The 17-case pilot design is deterministic and auditable. \\
\bottomrule
\end{longtable}

Every test calls \texttt{setup\_project\_paths}, disables plots and MAT-file saving, calls the numerical solver directly, and uses MATLAB \texttt{assert} statements. A printed \texttt{passed} line means all assertions in that function completed.

\subsection{Batch-test command}

\begin{lstlisting}[style=projectshell]
matlab -batch "cd('C:/path/to/RadiationEffects_proj2/matlab'); setup_project_paths; test_module2_fem_all_2d;"
\end{lstlisting}

\subsection{Mesh-refinement check}

Passing one fixed-mesh test is necessary but not sufficient for a research calculation. For a convergence study, repeat the same physical case on successively finer meshes and compare a scalar or field quantity. The uniform-charge case has an analytical nodal reference:
\begin{lstlisting}[style=projectmatlab]
meshSizes = [21, 41, 81, 161];
errors = zeros(size(meshSizes));

for k = 1:numel(meshSizes)
    params = default_module2_params('uniform_space_charge');
    params.nx = meshSizes(k);
    params.ny = 1 + (meshSizes(k)-1)/2;
    r = solve_poisson_defect_space_charge_2d(params);

    x = r.mesh.nodes(:,1);
    exact = (params.rho_uniform/(2*params.eps_si)) .* ...
        x .* (params.Lx-x);
    errors(k) = max(abs(r.phi-exact)) / max(max(abs(exact)), eps);
end

loglog(meshSizes, errors, 'o-');
xlabel('Number of x nodes'); ylabel('Relative max error'); grid on;
\end{lstlisting}
For this special problem, the exact solution is quadratic in $x$, whereas the T3 approximation is piecewise linear. Error should decrease with refinement. A refinement study for a coupled, highly localized source should additionally monitor peak potential, field magnitude, integrated charge, and spatial profiles.

\section{Debugging and interactive code tracing}

\subsection{Stop at the first MATLAB error}

\begin{lstlisting}[style=projectmatlab]
dbstop if error
result = main_module2_electrostatics('localized_defect_charge');
\end{lstlisting}
MATLAB will pause at the line that raised the error. Use the Workspace, \texttt{whos}, and the editor's step controls to inspect local variables. Clear the diagnostic stop later with \texttt{dbclear if error}.

\subsection{Enter the solver at its first line}

\begin{lstlisting}[style=projectmatlab]
dbstop in solve_poisson_defect_space_charge_2d at 1
result = main_module2_electrostatics('localized_defect_charge');
\end{lstlisting}
Step through in this order: \texttt{mesh}, \texttt{rho}, \texttt{K/rhs}, constrained nodes, \texttt{Kbc/rhsbc}, \texttt{phi}, then \texttt{field}. This sequence mirrors the actual numerical dependency chain.

\subsection{Inspect source code and resolution}

\begin{lstlisting}[style=projectmatlab]
open main_module2_electrostatics
open solve_poisson_defect_space_charge_2d
open assemble_poisson_fem_2d
which apply_dirichlet_bc -all
\end{lstlisting}
If edits are not reflected in execution, check \texttt{which -all}, then use \texttt{clear functions} and \texttt{rehash path}. Do not continue until MATLAB resolves the intended project copy.

\section{Common failure modes and diagnostic actions}

\begin{longtable}{p{0.28\textwidth} p{0.29\textwidth} p{0.35\textwidth}}
\toprule
\textbf{Symptom} & \textbf{Likely cause} & \textbf{Action} \\
\midrule
\endfirsthead
\toprule
\textbf{Symptom} & \textbf{Likely cause} & \textbf{Action} \\
\midrule
\endhead
\texttt{Undefined function} for a Module 2 file & Project paths were not added, or MATLAB is in the wrong extracted copy. & Run \texttt{setup\_project\_paths}; inspect \texttt{which ... -all}. \\
\texttt{Unknown Module 2 case name} & The driver accepts only the six strings in the default-parameter switch. & Use an exact built-in name or call the solver with an overridden parameter structure. \\
\texttt{Triangle has non-positive area} & Connectivity was changed to clockwise order, or the mesh contains a degenerate element. & Inspect the reported element coordinates and restore counter-clockwise connectivity. \\
Singular or nearly singular matrix warning & No Dirichlet reference potential, disconnected mesh, or invalid constraints. & Confirm \texttt{fixedNodes} is nonempty and the mesh is connected. \\
Changing \texttt{dphidn} has no effect & Nonzero Neumann boundary integration is not implemented. & Keep natural flux zero or extend assembly with the correct boundary integral. \\
Changing defect parameters has no effect & \texttt{rho\_uniform} exists and bypasses the general source calculation. & Remove the field or start from a non-uniform-source case. \\
NaN or Inf in charge & Zero/invalid Gaussian width or nonfinite input concentration. & Require positive \texttt{Cdef\_sigma\_x/y} and inspect every source field. \\
Unexpected potential sign & Charge sign, effective defect charge, or field convention was misread. & Inspect \texttt{rho}; remember positive source with grounded contacts gives a positive potential, while $\Evec=-\nabla\phi$. \\
Unexpected corner potential & Two Dirichlet sources prescribe different values at one node. & Make the values compatible; the collector now rejects conflicting constraints explicitly. \\
No output files after a direct solver call & Saving, plotting, and summary writing are driver responsibilities. & Call the three output functions explicitly or use the top-level driver. \\
Plot creation fails in a headless environment & Graphics renderer is unavailable even for invisible figures. & Use the direct solver route and omit \texttt{plot\_module2\_result\_2d}; inspect numerical arrays and summaries. \\
Excessive memory or long solve time & \texttt{nx} and \texttt{ny} create too many nodes/elements for the environment. & Start with the defaults, refine gradually, and monitor \texttt{whos result}. \\
Unphysical magnitude & A concentration or length was entered in centimeters, micrometers, or cubic centimeters without SI conversion. & Convert lengths to meters and concentrations to m$^{-3}$ before solving. \\
\bottomrule
\end{longtable}

\section{Adding a new standalone case safely}

There are two useful levels of customization.

\subsection{One-off experiment}

Do not edit the parameter factory. Create and override a structure in a script, call the solver, inspect the result, and save it under a unique name. This is the recommended exploratory path because it cannot silently alter the six built-in cases.

\subsection{New named regression or demonstration case}

For a permanent named case:
\begin{enumerate}[leftmargin=*]
  \item Add a new \texttt{case} branch in \texttt{default\_module2\_params.m}.
  \item Set every field that differs from the common defaults.
  \item Optionally add a convenience wrapper in \matlabfile{matlab/cases/}.
  \item Add a test in \matlabfile{matlab/tests/} with a quantitative pass criterion.
  \item Run \texttt{test\_module2\_fem\_all\_2d} to detect legacy or transmon regressions.
  \item Run the new case through the top-level driver and inspect every output.
  \item Document the case name, units, reference solution or expected trend, and limitations.
\end{enumerate}
Adding only a wrapper in \matlabfile{matlab/cases/} does not make the name acceptable to the current top-level driver; the switch in \texttt{default\_module2\_params.m} remains authoritative.

\section{Preparing for Module 3 coupling}

The standalone source builder currently constructs
\begin{equation}
\rho_{\mathrm{def}}(x,y)=qz_{\mathrm{def}}C_{\mathrm{def}}(x,y)
\end{equation}
from an internal Gaussian. In the coupled architecture, Module 3 should instead supply one or more nodal concentration fields $C_i(x,y,t)$ and Module 2 should evaluate
\begin{equation}
\rho_{\mathrm{def}}(x,y,t)=q\sum_i z_i C_i(x,y,t).
\end{equation}

Before implementing that interface, define and test:
\begin{itemize}[leftmargin=*]
  \item whether Modules 2 and 3 share exactly the same node ordering and mesh;
  \item the interpolation rule when their meshes differ;
  \item concentration units and effective charge signs;
  \item species ordering and metadata;
  \item conservation or interpolation error checks;
  \item behavior at times when no defect field is available;
  \item whether carrier and dopant contributions are supplied independently or already included in an imported total-charge map.
\end{itemize}

Do not overwrite \texttt{Cdef\_peak} with an array and assume the existing Gaussian expression will accept it correctly. Add an explicit nodal-field branch with dimension and unit validation, then protect it with a charge-consistency test.

\section{Recommended verification ladder before using a result}

Use the following order when changing the FEM code:
\begin{enumerate}[leftmargin=*]
  \item \textbf{Path check:} verify that MATLAB resolves the intended files.
  \item \textbf{Zero solution:} run \texttt{test\_module2\_zero\_charge\_2d}.
  \item \textbf{Boundary and sign check:} run \texttt{test\_module2\_linear\_potential\_2d}.
  \item \textbf{Source integration check:} run \texttt{test\_module2\_uniform\_space\_charge\_2d}.
  \item \textbf{Localized-source sanity:} run \texttt{test\_module2\_localized\_defect\_charge\_2d}.
  \item \textbf{Named-tag integration:} run \texttt{test\_module2\_transmon\_integration\_2d}.
  \item \textbf{Algebraic checks:} inspect residual, symmetry, fixed-node error, global reaction balance, and finiteness.
  \item \textbf{Mesh refinement:} establish convergence of the reported quantity.
  \item \textbf{Physical scale check:} compare charge, potential, and field magnitudes with an independent estimate.
  \item \textbf{Coupling check:} when importing fields, verify units, interpolation, sign, and total-charge consistency.
\end{enumerate}

\begin{physicsbox}
\textbf{Interpretation rule.} A successful sparse solve demonstrates that MATLAB found a vector satisfying the assembled algebraic equations. Analytical tests, convergence studies, dimensional checks, and coupling consistency are what establish that those equations represent the intended physical problem.
\end{physicsbox}

\section{Minimal handoff checklist}

Before sending or presenting a Module 2 FEM result, record:
\begin{itemize}[leftmargin=*]
  \item repository/archive version and MATLAB version;
  \item case name and any parameter overrides;
  \item geometry mode and, for transmon cases, Module 9 schema;
  \item $L_x$, $L_y$, \texttt{nx}, and \texttt{ny};
  \item permittivity and every nonzero charge contribution;
  \item boundary types, named tag sources, and values;
  \item passed tests and mesh-refinement evidence;
  \item algebraic residual, fixed-boundary error, and global reaction balance;
  \item output filenames and the exact quantity plotted;
  \item for batch data, parameter design version, complete-case split indices, operator/factorization counts, and worst per-case diagnostics;
  \item known limitations, especially prescribed sources and absent nonzero-Neumann support.
\end{itemize}

\section{Next development gate}

After the complete FEM suite passes, generate the 17-case pilot and inspect its text validation summary:
\begin{lstlisting}[style=projectmatlab]
[dataset, report] = main_module2_batch_fem_dataset();
\end{lstlisting}
The next coding gate is to adapt the Module 2 PINN to the Module 9 domain and this saved complete-case contract. The PINN must sample the substrate interior and named electrode segments, preserve the JJ as an interface rather than a short, normalize shifted $(x,z)$ coordinates and the five conditional parameters, and use the saved FEM potentials as reference/anchor data. The existing \matlabfile{module2_pinn_implementation_note.pdf} documents the current rectangular single-case PINN and must be revised when that conversion is implemented. Do not scale to hundreds of FEM configurations until this small end-to-end conditional PINN path passes on held-out pilot cases.

\end{document}
